\uuid{y3QF}
\exo7id{1675}
\titre{exo7 1675}
\auteur{ortiz}
\organisation{exo7}
\datecreate{1999-04-01}
\isIndication{false}
\isCorrection{false}
\chapitre{Réduction d'endomorphisme, polynôme annulateur}
\sousChapitre{Sous-espace stable}
\cpge{mpsi-mp2i.applications-lineaires.b, pcsi.applications-lineaires.b, ptsi.applications-lineaires.b, mpsi-mp2i.matrices-applications-lineaires.a, pcsi.matrices-applications-lineaires.a, ptsi.matrices-applications-lineaires.a, ptsi.geometrie-espace.e, bcpst1.geometrie.d, mp-mpi.reduction-endomorphismes.b, pc.reduction.a, psi.reduction-endomorphismes.a, pt.reduction-endomorphismes.a}
\module{Algèbre}
\niveau{L2}
\difficulte{}

\contenu{
\question{
Soit l'endomorphisme $f$ de $\Rr^3$ canoniquement
associ\'e \`a la matrice
$M=\left(\begin{smallmatrix}
1&1&0\\-1&2&1\\1&0&1\end{smallmatrix}\right)$. Le
plan $P$ d'\'equation $y+z=0$ est-il stable par
$f$ ? La droite $\text{vect}\left\{(1,1,1)\right\}$
est-elle stable par $f$ ?
}
}
